%%% On the Computational Complexity of Optimal Yose Move in Go
%%% via Combinatorial Game Theory
%%%
%%% Type A paper: PSPACE-hardness lower bound + matched polynomial upper bound
%%% under standard Berlekamp--Wolfe shape promise.

\documentclass[11pt,a4paper]{amsart}

%% ---------- packages ----------
\usepackage[utf8]{inputenc}
\usepackage{amsmath,amssymb,amsthm,mathtools}
\usepackage[margin=2.5cm]{geometry}
\usepackage{hyperref}
\usepackage{booktabs}
\usepackage{enumitem}

%% ---------- theorem environments ----------
\newtheorem{theorem}{Theorem}[section]
\newtheorem*{theorem*}{Theorem}
\newtheorem{lemma}[theorem]{Lemma}
\newtheorem{proposition}[theorem]{Proposition}
\newtheorem{corollary}[theorem]{Corollary}
\theoremstyle{definition}
\newtheorem{definition}[theorem]{Definition}
\theoremstyle{remark}
\newtheorem{remark}[theorem]{Remark}
\newtheorem{example}[theorem]{Example}

%% ---------- macros ----------
\newcommand{\R}{\mathbb{R}}
\newcommand{\Q}{\mathbb{Q}}
\newcommand{\Z}{\mathbb{Z}}
\newcommand{\N}{\mathbb{N}}
\newcommand{\LS}{\mathrm{LS}}
\newcommand{\RS}{\mathrm{RS}}
\newcommand{\Tt}{\mathrm{T}}
\newcommand{\PSPACE}{\mathsf{PSPACE}}
\newcommand{\GLG}{\mathsf{GLG}}
\newcommand{\BMX}{\textsc{Best-Move-X}}

%% ---------- metadata ----------
\title[Complexity of Optimal Yose Move in Go]{On the Computational Complexity\\
  of the Optimal Yose Move in Go via Combinatorial Game Theory}

\author{Lightman Chang}
\address{Independent Researcher}
\email{lightman.chang@gmail.com}

\date{\today}

\subjclass[2020]{Primary 91A46; Secondary 68Q17, 91A05}
\keywords{Combinatorial game theory, Go endgame, computational complexity,
  PSPACE-hardness, thermograph, Berlekamp--Wolfe theory}

\begin{document}
\maketitle

%% =================================================================
\begin{abstract}
We study the complexity of computing an optimal endgame (\emph{yose}) move
in $n \times n$ Go under three formal readings of ``optimal'': maximum
combinatorial game theory (CGT) temperature ($X = B$); maximum
swing $\LS - \RS$ ($X = C$); and maximum (or minimum) value-of-being-forced-to-play
under optimal sequel ($X = D\text{-must}$, $X = D\text{-avoid}$). Working
within the local decomposition framework of Berlekamp and Wolfe, we prove
that for every $X \in \{B, C, D\text{-must}, D\text{-avoid}\}$ the problem
$\BMX$ is $\PSPACE$-hard, by a reduction from the Generalized Ladder Game
(Cr\^{a}\c{s}maru--Tromp 2000). The reduction is calibrated by an explicit
\emph{swing-booster gadget} (Lemma~\ref{lem:swing}) that yields a polynomial
gap $D_Y - D_N = \Omega(|B|)$ between yes-instances and no-instances. We
complement the lower bound with a matched algorithmic upper bound: under the
standard Berlekamp--Wolfe shape promise, ko-freeness, and local
decomposability, $\BMX$ admits an
$O(n^2 + m \log m)$ algorithm based on closed-form region values and an
$O(m)$ prefix-sum evaluation of the optimal-sequel scores. The two
results together identify the precise barrier between tractable and
intractable yose, and clarify which formalizations of ``best move'' are
algorithmically equivalent.
\end{abstract}

%% =================================================================
\section{Introduction}\label{sec:intro}
%% =================================================================

The closing phase of a game of Go---the \emph{yose}---admits a natural
formalization within combinatorial game theory (CGT), where the board
decomposes into disjoint local games whose values are added under the
disjunctive sum operation \cite{Conway1976,BCG1982}. In the classical
Berlekamp--Wolfe theory of cooled Go positions \cite{BW1994}, every
``standard'' shape (corridor, simple switch, sente sequence, etc.) admits
explicit polynomial-time stops and temperature, and the player who maximizes
the global score is the player who maximizes a known function of these
local values. This raises a basic complexity question:

\begin{quote}
\emph{Given a Go position in the yose phase, how hard is it to compute the
local move that is optimal under a precisely specified definition of
``best''?}
\end{quote}

Three definitions of ``best move'' are standard in the CGT and Go
literature, and they need not coincide:
\begin{itemize}[leftmargin=2em]
\item[$(B)$] \emph{Hottest move}: the local game $R_i$ with the largest
  CGT temperature $\Tt(R_i)$.
\item[$(C)$] \emph{Largest swing}: the local game with the largest swing
  $\Delta(R_i) = \LS(R_i) - \RS(R_i)$.
\item[$(D)$] \emph{Win-deciding move}: the local game whose forced
  selection most affects the final score under optimal subsequent play.
\end{itemize}

Definition $(D)$ admits at least three precise readings (must-play,
avoid-play, sign-flipping locus); see Definition~\ref{def:Dreadings}.
Throughout the paper we restrict attention to the must-play and avoid-play
readings; the sign-flipping reading involves a self-referential baseline
$V_0$ and is left open.

\paragraph{Background.}
Robson \cite{Robson1983} proved that $n \times n$ Go is
$\mathsf{EXPTIME}$-complete using nested ko constructions. This bound is
\emph{global}: the entire board, with all the complications of ko and
life-and-death, is encoded into a single instance. For the
\emph{yose} phase, ko has typically been resolved and the position
admits a CGT decomposition. In this restricted regime,
Cr\^{a}\c{s}maru and Tromp~\cite{CrasmaruTromp2000} showed that even the
single-region problem of deciding the outcome of a Go ladder is
$\PSPACE$-complete, by encoding QBF into the ladder dynamics. Their
construction is a natural building block for a complexity theory of the
\emph{yose} that lies strictly below the $\mathsf{EXPTIME}$ bound of
Robson.

On the upper-bound side, Berlekamp and Wolfe~\cite{BW1994} catalogued a
family of standard shapes whose CGT canonical form is computable by closed
formulas. Modulo the catalogue and the disjunctive-sum law of Conway, this
yields polynomial-time algorithms for $(B)$ and $(C)$ on positions
composed entirely of standard shapes.

\paragraph{The gap.}
Two questions remain unresolved in this picture. First, under the
\emph{generic} local-decomposition promise (no shape restriction), is
computing the optimal yose move actually as hard as deciding QBF? Second,
do the three definitions $(B)$, $(C)$, $(D)$ produce the same
top-$1$ choice on standard shapes, and if not, where do they diverge?

\paragraph{Our contribution.}
We address the first question for all four readings
$X \in \{B, C, D\text{-must}, D\text{-avoid}\}$. Our main result
(Theorem~\ref{thm:main}) is:

\begin{theorem*}[Informal]
For every $X \in \{B, C, D\text{-must}, D\text{-avoid}\}$, the problem
$\BMX$ is $\PSPACE$-hard, even under the local-decomposition promise.
Conversely, under the additional Berlekamp--Wolfe standard-shape and
ko-free promises, $\BMX$ admits an $O(n^2 + m \log m)$ algorithm,
where $n$ is the board side and $m$ is the number of regions.
\end{theorem*}

The key technical step in the lower bound is a \emph{swing-control lemma}
(Lemma~\ref{lem:swing}): given any Generalized Ladder Game instance
$(B, s_0)$, we construct in polynomial time a region $R_\Phi$ with an
explicit \emph{escape booster} and a \emph{capture booster} gadget, such that
\[
  \GLG(B, s_0) = \mathrm{YES}
  \Longleftrightarrow
  \Delta(R_\Phi) = D_Y \geq 6k,
  \qquad
  \GLG(B, s_0) = \mathrm{NO}
  \Longleftrightarrow
  \Delta(R_\Phi) = D_N \leq 4,
\]
where $k = \mathrm{poly}(|B|)$ is the chase length of the ladder. The
polynomial gap $D_Y - D_N = \Omega(|B|)$ closes a calibration gap left in
prior informal arguments and allows us to instantiate a single reduction
template for all four definitions of ``best move''.

The matching upper bound (Theorem~\ref{thm:upper}) reduces optimal-sequel
scores to a closed-form prefix-sum expression over the swing-sorted
permutation of regions, giving $O(m)$ work per region after sorting.

\paragraph{Comparison with prior work.}
Robson's $\mathsf{EXPTIME}$-completeness theorem~\cite{Robson1983}
operates on the unrestricted Go-position problem and uses ko, while our
lower bound holds in the ko-free yose regime---hence our bound is
strictly weaker ($\PSPACE$ vs. $\mathsf{EXPTIME}$) but applies to a
strictly more restricted class of inputs. The Cr\^{a}\c{s}maru--Tromp
result~\cite{CrasmaruTromp2000} established $\PSPACE$-completeness for a
single ladder, but did not address the comparison-of-regions structure
that distinguishes ``best move'' problems from outcome problems. Our
swing-booster construction (Lemma~\ref{lem:swing}) and the dominance
analysis of its CGT canonical form (Lemma~\ref{lem:canonical}) are new.

\paragraph{Paper structure.}
Section~\ref{sec:prelim} fixes CGT notation and defines the four
problems $\BMX$. Section~\ref{sec:lower} states and proves the swing-control
lemma and the $\PSPACE$-hardness theorem. Section~\ref{sec:upper}
gives the polynomial-time algorithm under the standard-shape promise.
Section~\ref{sec:discussion} discusses the gap between the two bounds and
open problems.

%% =================================================================
\section{Preliminaries}\label{sec:prelim}
%% =================================================================

\subsection{Combinatorial game values}

We work throughout in the framework of Conway~\cite{Conway1976} and
Berlekamp--Conway--Guy~\cite{BCG1982}. A \emph{game} $G$ is recursively
defined by a set of \emph{Left options} $G^L$ and \emph{Right options} $G^R$.
For a game $G$, the \emph{Left stop} and \emph{Right stop} are
\[
  \LS(G) = \max_{G^L \in G^L} \RS(G^L),
  \qquad
  \RS(G) = \min_{G^R \in G^R} \LS(G^R),
\]
with the base case $\LS(G) = \RS(G) = G$ when $G$ is a (dyadic rational)
number. The \emph{swing} of $G$ is $\Delta(G) = \LS(G) - \RS(G) \geq 0$,
its \emph{mean} is $\mu(G) = (\LS(G) + \RS(G))/2$, and its
\emph{temperature} $\Tt(G)$ is the height at which the two scaffolds of the
thermograph of $G$ meet \cite[Ch.~6]{BCG1982}. For a simple switch
$S = \{a \mid b\}$ with $a > b$ both numbers, $\Tt(S) = (a-b)/2$ and
$\mu(S) = (a+b)/2$.

\subsection{Local decomposition and the yose phase}

\begin{definition}[Local decomposability]\label{def:decomp}
A Go position $s$ on an $n \times n$ board is \emph{locally decomposable}
if its set of empty intersections partitions into disjoint regions
$\{R_1, \ldots, R_m\}$ such that
\begin{enumerate}[label=(\roman*)]
\item for every legal play sequence, the liberties and stones inside
  $R_i$ depend only on play inside $R_i$;
\item the global Japanese-rules score satisfies
  $V(s) = \sum_{i=1}^m V_i(R_i)$;
\item each $R_i$ defines a well-formed CGT subgame.
\end{enumerate}
\end{definition}

\begin{definition}[Ko-free promise]\label{def:kofree}
A region $R_i$ is \emph{ko-free} if its local game tree contains no node
whose subboard (restricted to $R_i$ together with its boundary) recurs at
any descendant. This is strictly stronger than the simple ko rule of
Japanese rules.
\end{definition}

We assume throughout that the input is locally decomposable (a promise).
The yose phase is, by definition, the regime in which this promise holds.

\subsection{Definitions of optimal move}

\begin{definition}[Three readings of ``best
move'']\label{def:Dreadings}
Let $s$ have decomposition $\{R_i\}_{i=1}^m$. Let
\[
  V_j \;=\; \RS\Bigl( R_j^{L^*} + \sum_{i \neq j} R_i \Bigr),
  \quad\text{where } R_j^{L^*}
  = \arg\max_{R_j^L} \RS\Bigl( R_j^L + \sum_{i \neq j} R_i \Bigr).
\]
That is, $V_j$ is the final score after Black is forced to play in $R_j$
and both sides play optimally thereafter. We define:
\begin{itemize}[leftmargin=2em]
\item[$(D\text{-must})$] $j^* = \arg\max_j V_j$. ``The move Black should
  play to maximize the final score.''
\item[$(D\text{-avoid})$] $j^* = \arg\min_j V_j$. ``The move Black must
  not play.''
\item[$(D\text{-sign})$] Order regions by $|V_j - V_0|$, where
  $V_0 = \LS(\sum_i R_i)$ is the unconstrained optimal-play value.
\end{itemize}
\end{definition}

\begin{definition}[The problems $\BMX$]\label{def:BMX}
For $X \in \{B, C, D\text{-must}, D\text{-avoid}\}$, the problem
$\BMX$ takes as input a locally decomposable position $s$ and an integer
$k \geq 1$, and outputs the top-$k$ regions ranked by
\[
  X = B: \;\Tt(R_i) \text{ descending};
  \qquad
  X = C: \;\Delta(R_i) \text{ descending};
\]
\[
  X = D\text{-must}: \;V_i \text{ descending};
  \qquad
  X = D\text{-avoid}: \;V_i \text{ ascending}.
\]
\end{definition}

The sign-flipping reading $(D\text{-sign})$ requires the baseline $V_0$,
which is itself the answer to a $\BMX$-style optimization, leading to
a circular dependence we do not resolve here; see
Section~\ref{sec:discussion}.

\subsection{Crâșmaru--Tromp Generalized Ladder Game}

We rely on the following theorem.

\begin{theorem}[Cr\^{a}\c{s}maru--Tromp \cite{CrasmaruTromp2000}]
\label{thm:GLG}
The Generalized Ladder Game ($\GLG$) is $\PSPACE$-complete: given a
rectangular Go region $B$ (enclosed by a 2-stone-wide live wall) with a
chaser stone, an escaper stone, and obstruction stones at fixed
positions, deciding whether the escaper has a winning strategy under the
standard ladder dynamics is $\PSPACE$-complete.
\end{theorem}

The proof reduces from QBF, the prototypical $\PSPACE$-complete problem
\cite{StockmeyerMeyer1973}: each quantifier alternation is encoded as a
turn in the ladder where the current player's choice corresponds to a
variable assignment; obstructions enforce the ``wrong choice $\Rightarrow$
forced capture'' structure. The complete construction occupies the
30 pages of \cite{CrasmaruTromp2000}; we use it as a black box. The
overall reduction structure is in the spirit of the polynomial-time
many-one reductions of \cite{Karp1972}.

%% =================================================================
\section{Lower Bound: PSPACE-Hardness}\label{sec:lower}
%% =================================================================

The main lower bound is the following theorem; the heart of the argument
is the swing-control lemma below.

\begin{theorem}\label{thm:main}
For every $X \in \{B, C, D\text{-must}, D\text{-avoid}\}$, the problem
$\BMX$ is $\PSPACE$-hard, even under the local-decomposition promise.
\end{theorem}

\subsection{Swing-control via booster gadgets}

\begin{lemma}[Swing-control lemma]\label{lem:swing}
There is a polynomial-time algorithm that, given any $\GLG$ instance
$(B, s_0)$ with chase length $k = \mathrm{poly}(|B|)$, constructs a Go
region $R_\Phi$ embedded inside a 2-stone-wide live white wall, such that
\begin{align*}
  \GLG(B, s_0) = \mathrm{YES}
   &\;\Longrightarrow\; \LS(R_\Phi) = 4k,\;
   \RS(R_\Phi) = -2k, \;\Delta(R_\Phi) =: D_Y = 6k,
   \\
  \GLG(B, s_0) = \mathrm{NO}
   &\;\Longrightarrow\; \LS(R_\Phi) = -2k + \varepsilon,\;
   \RS(R_\Phi) = -2k - \varepsilon',
   \;\Delta(R_\Phi) =: D_N = \varepsilon + \varepsilon' \leq 4,
\end{align*}
where $\varepsilon, \varepsilon' \in [0, 2]$. In particular,
$D_Y - D_N \geq 6k - 4 = \Omega(|B|)$.
\end{lemma}

\begin{proof}
We give an explicit construction.
The region $R_\Phi$ embeds the ladder block $B$ with its obstruction
stones, plus two booster gadgets attached at the two distinguished
endpoints of the ladder dynamics:
\begin{itemize}[leftmargin=2em]
\item \emph{Escape booster}: a $1 \times 2k$ corridor attached at the
  ``successful escape'' terminal of the ladder. If the escaper succeeds,
  Black additionally collects $2k$ points by occupying the corridor in a
  forced sente sequence (standard corridor evaluation,
  cf.~\cite[Ch.~3]{BW1994}).
\item \emph{Capture booster}: a small dead shape (e.g., a 4-stone
  ``one-eye'' configuration) attached at the ``failed escape'' terminal.
  Both colors play to no effect inside this gadget, contributing
  $\LS = \RS = 0$.
\end{itemize}
The total size of the boosters is $O(k)$, which fits inside the
2-stone-wide white wall enclosure of $R_\Phi$ for sufficiently large
$n$, and remains decoupled from the rest of the board.

\medskip
\noindent\textbf{Computation of $\LS(R_\Phi)$ and $\RS(R_\Phi)$.}
Let $k$ be the chase length of the ladder.

\emph{Case $\GLG = \mathrm{YES}$.}
If Black plays first in $R_\Phi$, she initiates the chase, escapes
successfully (saving $k$ stones, $+2k$ points), and then collects the
escape booster ($+2k$ points), giving $\LS(R_\Phi) = 4k$.
If White plays first, the chase forces the capture of Black's $k$ stones
($-2k$ points) and the boosters are inert, giving $\RS(R_\Phi) = -2k$.

\emph{Case $\GLG = \mathrm{NO}$.}
If Black plays first, she still attempts the chase but the escape fails;
the forced sequence captures her $k$ stones, with a small move-parity
correction $\varepsilon \in [0, 2]$ from the boundary of the chase. So
$\LS(R_\Phi) = -2k + \varepsilon$. Symmetrically,
$\RS(R_\Phi) = -2k - \varepsilon'$ with $\varepsilon' \in [0, 2]$.

The values $\varepsilon, \varepsilon'$ are fully determined by the
ladder dynamics on the boundary stones and can be computed in polynomial
time from $(B, s_0)$. In particular,
$D_N = \varepsilon + \varepsilon' \leq 4$.
\end{proof}

\begin{lemma}[CGT canonical form of $R_\Phi$]\label{lem:canonical}
Under the construction of Lemma~\ref{lem:swing}, the canonical form of
$R_\Phi$ in the sense of \cite[Ch.~4]{BCG1982} is the simple switch
$\{\LS(R_\Phi) \mid \RS(R_\Phi)\}$.
\end{lemma}

\begin{proof}
We show that all Left options of $R_\Phi$ except the ``initiate-chase''
option are dominated, and similarly for Right options.

The Left options of $R_\Phi$ are:
\begin{itemize}[leftmargin=2em]
\item[$(L_1)$] Initiate the ladder chase. Final value:
  $4k$ (YES case), $-2k + \varepsilon$ (NO case).
\item[$(L_2)$] Play directly inside the escape booster corridor.
  The corridor is sealed by a live white wall at this point in the game
  tree, so a Black stone there has no liberties and is immediately
  captured. The chase still proceeds and is forced by White. Final value
  $\leq -2k - 1$.
\item[$(L_3)$] Play directly inside the capture-booster dead shape. The
  shape is dead, so the stone is wasted, and White still drives the
  chase. Final value $= -2k$.
\end{itemize}

In the YES case, $(L_1)$ strictly dominates: $4k$ versus $\leq -2k - 1$
and $-2k$. In the NO case, $(L_1)$ achieves $-2k + \varepsilon \geq -2k$,
which weakly dominates $(L_3)$ and strictly dominates $(L_2)$. The
canonical-form algorithm of \cite[Ch.~4]{BCG1982} eliminates dominated
options under the weak inequality $\geq$, so $(L_2), (L_3)$ are removed
and only $(L_1)$ remains.

By symmetry (the Right options $(R_1)$, $(R_2)$, $(R_3)$ are analogous),
only the symmetric ``initiate-chase'' option $(R_1)$ survives among Right
options. Each surviving option's value is a number, so the canonical
form is the simple switch
$\{\LS(R_\Phi) \mid \RS(R_\Phi)\}$.
\end{proof}

\subsection{Reduction template}

Fix a $\GLG$ instance $(B, s_0)$ and apply Lemma~\ref{lem:swing} to obtain
$R_\Phi$ with parameters $D_Y, D_N$. Let $a := (D_Y + D_N) / 4$, which is
computable in polynomial time, and define a second region
$R_*$ to be a simple switch $\{a \mid -a\}$ (realizable in Go as a
standard $1 \times \ell$ corridor with $\ell$ chosen so that
$\Delta(R_*) = 2a$; see~\cite[Ch.~3]{BW1994} for the catalogue). Both
$R_\Phi$ and $R_*$ sit inside disjoint, sealed sub-boards, so the
two-region position is locally decomposable in the sense of
Definition~\ref{def:decomp}.

\begin{proof}[Proof of Theorem~\ref{thm:main}]
We show that for each $X$, an oracle for $\BMX$ on the two-region
instance $\{R_\Phi, R_*\}$ decides $\GLG$.

\emph{Case $X = B$ (temperature).} Since $R_\Phi$ has canonical form
$\{\LS(R_\Phi) \mid \RS(R_\Phi)\}$ by Lemma~\ref{lem:canonical}, we have
$\Tt(R_\Phi) = \Delta(R_\Phi) / 2$. Also $\Tt(R_*) = a$. Hence
\[
  \Tt(R_\Phi) > \Tt(R_*)
  \;\Longleftrightarrow\;
  \Delta(R_\Phi) > 2a
  \;\Longleftrightarrow\;
  D_Y > 2a > D_N
  \;\Longleftrightarrow\;
  \GLG = \mathrm{YES}.
\]
The middle equivalence uses $a = (D_Y + D_N)/4$, which lies strictly
between $D_N/2$ and $D_Y/2$ since $D_Y > D_N$ (Lemma~\ref{lem:swing}).

\emph{Case $X = C$ (swing).} Identical to $X = B$ since
$\Delta(R_\Phi) > \Delta(R_*) = 2a$ iff $\GLG = \mathrm{YES}$.

\emph{Case $X = D\text{-must}$.} For the two-region instance,
\begin{align*}
  V_{R_\Phi} &= \LS(R_\Phi) + \RS(R_*) = \LS(R_\Phi) - a,
  \\
  V_{R_*} &= \LS(R_*) + \RS(R_\Phi) = a + \RS(R_\Phi).
\end{align*}
Hence
\[
  V_{R_\Phi} - V_{R_*}
  = \bigl(\LS(R_\Phi) - \RS(R_\Phi)\bigr) - 2a
  = \Delta(R_\Phi) - 2a,
\]
which is positive iff $\GLG = \mathrm{YES}$.

\emph{Case $X = D\text{-avoid}$.} The same calculation gives
$V_{R_\Phi} < V_{R_*} \iff \Delta(R_\Phi) < 2a \iff \GLG = \mathrm{NO}$,
so the avoid-play top-$1$ identifies $R_\Phi$ iff $\GLG = \mathrm{NO}$.

In all four cases the reduction is polynomial in $|B|$ and decides
$\GLG$. By Theorem~\ref{thm:GLG}, $\BMX$ is $\PSPACE$-hard.
\end{proof}

%% =================================================================
\section{Upper Bound under the Standard-Shape Promise}\label{sec:upper}
%% =================================================================

We now show that the lower bound is essentially tight: under a
shape-restriction promise that excludes the ladder constructions of
Section~\ref{sec:lower}, $\BMX$ becomes polynomial.

\begin{definition}[Berlekamp--Wolfe standard shape]\label{def:BWshape}
A region $R$ is a \emph{Berlekamp--Wolfe standard shape} if its CGT
canonical form is a finite disjunctive sum of the following components:
\begin{enumerate}[label=(\alph*)]
\item a number (a region with no hot move);
\item a simple switch $\{a \mid b\}$ with $a, b$ numbers and $a > b$;
\item a corridor of length $\ell$ with standard live-stone closure
  \cite[Ch.~3]{BW1994};
\item a finite sente sequence terminating in (a) or (b);
\item a dyadic infinitesimal in $\{0, *, \uparrow, \downarrow, \uparrow*,
  \cdot 2, \dots\}$.
\end{enumerate}
\end{definition}

For each component type, closed-form values $(\LS, \RS, \Tt, \mu)$ and an
optimal move are computable in $O(|R|)$ by standard formulas; the
infinitesimal components contribute $\LS = \RS = 0$ and do not affect
top-$k$ rankings of hot regions. We refer to the full procedure as
$\textsc{BWFormula}(R)$.

\subsection{The hotstrat-on-switches lemma}

A key combinatorial step is the following identity, which expresses the
optimal-sequel value of a sum of simple switches in closed form.

\begin{lemma}[Hotstrat on simple switches]\label{lem:H}
Let $G = \sum_{i=1}^m S_i$ with $S_i = \{a_i \mid b_i\}$ simple switches
and $\Delta_i = a_i - b_i > 0$. Sort the indices in descending order of
$\Delta$: $\Delta_{\sigma(1)} \geq \cdots \geq \Delta_{\sigma(m)}$. Then
\begin{enumerate}[label=(\alph*)]
\item under optimal play with Black to move, Black's first move is in
  $S_{\sigma(1)}$;
\item the entire optimal play alternates by selecting the
  largest remaining $\Delta$, with Black taking $a_{\sigma(l)}$ at odd
  positions $l$ and White taking $b_{\sigma(l)}$ at even positions;
\item $\LS(G) = \sum_{l \text{ odd}} a_{\sigma(l)} +
  \sum_{l \text{ even}} b_{\sigma(l)}$, and the symmetric formula holds
  for $\RS(G)$.
\end{enumerate}
\end{lemma}

\begin{proof}
We argue by an adjacent-swap exchange.
View the play as a sequence $\pi : [m] \to [m]$ that selects the
$\pi(l)$-th switch at position $l$, with Black at odd positions and
White at even positions. The score is
$\mathrm{Score}(\pi) = \sum_{l \text{ odd}} a_{\pi(l)} +
  \sum_{l \text{ even}} b_{\pi(l)}$.
For any $l \in [1, m-1]$, swapping positions $l$ and $l+1$ yields
\[
  \mathrm{Score}(\pi) - \mathrm{Score}(\pi')
  = \pm\bigl(\Delta_{\pi(l)} - \Delta_{\pi(l+1)}\bigr),
\]
with sign $+$ when $l$ is odd (Black's choice) and $-$ when $l$ is even
(White's choice). At a saddle point the controlling player gains nothing
by swapping, which forces $\Delta_{\pi^*(l)} \geq \Delta_{\pi^*(l+1)}$ in
both cases. Hence $\pi^* = \sigma$, and the formulas in (a)--(c) follow
directly.
\end{proof}

\subsection{The polynomial-time algorithm}

\begin{theorem}\label{thm:upper}
Suppose the input position $s$ on an $n \times n$ board is locally
decomposable, ko-free, and every region is a Berlekamp--Wolfe standard
shape. Then for every
$X \in \{B, C, D\text{-must}, D\text{-avoid}\}$ the problem $\BMX$ is
solvable in time $O(n^2 + m \log m)$.
\end{theorem}

\begin{proof}
We give algorithms for each $X$.

\smallskip\noindent\emph{Algorithm for $X = B$.}
\begin{enumerate}[label=\arabic*.]
\item Decompose $s$ into $\{R_1, \ldots, R_m\}$ via BFS in $O(n^2)$.
\item For each $R_i$, compute $(\LS_i, \RS_i, \Tt_i, m_i^*)
  \leftarrow \textsc{BWFormula}(R_i)$ in $O(|R_i|)$ total $O(n^2)$.
\item Sort regions by $\Tt_i$ in descending order; output top $k$.
\end{enumerate}
The sort costs $O(m \log m)$. Total: $O(n^2 + m \log m)$.

\smallskip\noindent\emph{Algorithm for $X = C$.}
Same as $X = B$ but sort by $\Delta_i = \LS_i - \RS_i$.

\smallskip\noindent\emph{Algorithm for $X = D\text{-must}$.}
We use Lemma~\ref{lem:H}. Let $\sigma$ be the descending order of
$\Delta$ over the hot regions, and write $a_l = \LS_{\sigma(l)}$,
$b_l = \RS_{\sigma(l)}$. By the lemma applied to ``Black is forced to
play in region $\sigma(p)$ first'', the optimal-sequel total is
\[
  V_{\sigma(p)} = a_p
  + \!\!\sum_{\substack{l \leq p-1 \\ l \text{ even}}}\!\! a_l
  + \!\!\sum_{\substack{l \leq p-1 \\ l \text{ odd}}}\!\! b_l
  + \!\!\sum_{\substack{l \geq p+1 \\ l \text{ odd}}}\!\! a_l
  + \!\!\sum_{\substack{l \geq p+1 \\ l \text{ even}}}\!\! b_l.
\]
Precompute the four prefix/suffix arrays
\[
  P_a^{\mathrm{ev}}[p] = \sum_{l \leq p,\, l \text{ even}} a_l,
  \;\;
  P_b^{\mathrm{od}}[p] = \sum_{l \leq p,\, l \text{ odd}} b_l,
  \;\;
  S_a^{\mathrm{od}}[p] = \sum_{l \geq p,\, l \text{ odd}} a_l,
  \;\;
  S_b^{\mathrm{ev}}[p] = \sum_{l \geq p,\, l \text{ even}} b_l,
\]
in $O(m)$ by single sweeps. Then each $V_{\sigma(p)}$ is computable in
$O(1)$ via
\[
  V_{\sigma(p)} = a_p + P_a^{\mathrm{ev}}[p-1] + P_b^{\mathrm{od}}[p-1]
    + S_a^{\mathrm{od}}[p+1] + S_b^{\mathrm{ev}}[p+1],
\]
giving $O(m)$ total over all regions, and a final $O(m \log m)$ sort.

\smallskip\noindent\emph{Algorithm for $X = D\text{-avoid}$.}
Same as $X = D\text{-must}$ but sort $V_{\sigma(p)}$ in ascending order.

In each case, the total time is $O(n^2 + m \log m)$.
\end{proof}

\subsection{Numerical sanity checks}

We verify the prefix-sum formula on small $m$.

\begin{example}\label{ex:m3}
Take $m = 3$ with $\Delta = (10, 8, 6)$ and $\mu_i = 0$, so
$(a_l, b_l) = (5, -5), (4, -4), (3, -3)$. Then
$P_a^{\mathrm{ev}} = [0,0,4,4]$, $P_b^{\mathrm{od}} = [0,-5,-5,-8]$,
$S_a^{\mathrm{od}} = [0,8,3,3,0]$ (indices $0$--$4$),
$S_b^{\mathrm{ev}} = [0,-4,-4,0,0]$. Direct evaluation gives
$V_1 = 5 + 0 + 0 + 3 - 4 = 4$, $V_2 = 4 + 0 - 5 + 3 + 0 = 2$,
$V_3 = 3 + 4 - 5 + 0 + 0 = 2$, matching the closed-form computation.
\end{example}

\begin{example}\label{ex:m4}
Take $m = 4$ with $\Delta = (10, 8, 6, 4)$ and $\mu_i = 0$. The same
formula gives $V_1 = 2$, $V_2 = V_3 = 0$, $V_4 = -2$, again matching the
recursive evaluation. Note the tie $V_2 = V_3$, foreshadowing the
structural-ties phenomenon studied in the companion paper.
\end{example}

%% =================================================================
\section{Discussion}\label{sec:discussion}
%% =================================================================

\subsection{The gap between the bounds}

Combining Theorems~\ref{thm:main} and~\ref{thm:upper}, we have a clean
dichotomy: the polynomial-time tractability of $\BMX$ is governed
precisely by the Berlekamp--Wolfe standard-shape promise. Without the
shape promise, the embedding of a ladder into a single region suffices to
make $\BMX$ as hard as deciding QBF.

An intermediate regime is also of interest. If we permit arbitrary
regions but require $|R_i| = O(\log n)$, then a brute-force traversal of
each local game tree gives an upper bound of
$n^{O((\log n)^2)}$---quasi-polynomial. The condition $|R_i| = O(\log n)$
thus appears to be the natural fine-grained boundary: above it
(constant-region-size, BW-standard) one is in $\mathsf{P}$, below it
(unrestricted, polynomial-region-size) one is $\PSPACE$-hard.

\subsection{The fourth reading: sign-flipping}

The sign-flipping reading $(D\text{-sign})$ orders regions by
$|V_j - V_0|$, where $V_0 = \LS(\sum_i R_i)$. The lower-bound reduction of
Section~\ref{sec:lower} does not directly apply, because the calibration
constant $a$ depends on the sign of $V_0$, which itself depends on the
$\GLG$ answer through the booster-corridor contribution. A careful
analysis would either eliminate this circularity or show that
$(D\text{-sign})$ is intrinsically harder than the must/avoid readings.
We leave this as an open question.

\subsection{Open problems}

\begin{enumerate}[label=\arabic*.]
\item \emph{$\mathsf{EXPTIME}$-hardness of yose under ko-freeness.}
  Robson's $\mathsf{EXPTIME}$-hardness theorem~\cite{Robson1983} uses
  global ko constructions. Is there a yose construction (in our sense)
  that is $\mathsf{EXPTIME}$-hard?
\item \emph{Multi-ko in the BW framework.} Extending the standard-shape
  catalogue to allow multi-ko situations is a long-standing open
  problem; even a partial polynomial-time algorithm under bounded
  ko-threat structure would be of independent interest.
\item \emph{Strong NP-hardness of top-$k$.} On general (non-standard)
  shapes, is the top-$k$ version of $\BMX$ NP-hard via an explicit
  $\textsc{Partition}$-style reduction? The closed-form ranking on
  standard shapes (Lemma~\ref{lem:H}) rules out the most direct attempt.
\item \emph{Empirical coverage.} What fraction of typical $19 \times 19$
  endgame positions falls into the BW standard-shape class? An empirical
  study would clarify the practical relevance of Theorem~\ref{thm:upper}.
\end{enumerate}

%% =================================================================
\begin{thebibliography}{10}

\bibitem{BCG1982}
E.~Berlekamp, J.~H. Conway, and R.~K. Guy.
\newblock \emph{Winning Ways for Your Mathematical Plays}.
\newblock Academic Press, 1982. (Second edition, A K Peters, 2001--2004.)

\bibitem{BW1994}
E.~Berlekamp and D.~Wolfe.
\newblock \emph{Mathematical Go: Chilling Gets the Last Point}.
\newblock A K Peters, 1994.

\bibitem{Conway1976}
J.~H. Conway.
\newblock \emph{On Numbers and Games}.
\newblock Academic Press, 1976. (Second edition, A K Peters, 2001.)

\bibitem{CrasmaruTromp2000}
M.~Cr\^{a}\c{s}maru and J.~Tromp.
\newblock Ladders are {PSPACE}-complete.
\newblock In \emph{Proceedings of the 2nd International Conference on Computers
  and Games (CG~2000)}, Lecture Notes in Computer Science, vol.~2063,
  pp.~241--249. Springer, 2001.

\bibitem{Robson1983}
J.~M. Robson.
\newblock The complexity of {G}o.
\newblock In \emph{Information Processing 83 (Proc. IFIP World Computer
  Congress)}, pp.~413--417. North-Holland, 1983.

\bibitem{StockmeyerMeyer1973}
L.~J. Stockmeyer and A.~R. Meyer.
\newblock Word problems requiring exponential time.
\newblock In \emph{Proceedings of the 5th Annual ACM Symposium on Theory of
  Computing (STOC)}, pp.~1--9. ACM, 1973.

\bibitem{Karp1972}
R.~M. Karp.
\newblock Reducibility among combinatorial problems.
\newblock In R.~E. Miller and J.~W. Thatcher, editors, \emph{Complexity of
  Computer Computations}, pp.~85--103. Plenum Press, 1972.

\end{thebibliography}

\end{document}
